\GalSourcePageBoundary{071}{L0070}{42}{42}

\GalHistoricalNote{tion en $r$ est de la forme $r^{p}=A$, et que les racines $p^{\text{ièmes}}$ de l'unité se trouvent
au nombre des quantités précédemment adjointes, les $p$ groupes dont il est
question dans le Théorème II jouiront, de plus, de cette \GalPrintedError{proprieté}{W09-PE002} que les
substitutions de lettres par lesquelles on passe d'une permutation à une
autre dans chaque groupe soient les mêmes pour tous les groupes. En effet,
dans ce cas, il revient au même d'adjoindre à l'équation telle ou telle valeur
de $r$. Par conséquent, ses propriétés doivent être les mêmes après l'adjonction de
telle ou telle valeur. Ainsi son groupe doit être le même quant aux substitutions
(Proposition I, scolie). Donc, etc.

Tout cela est effacé avec soin; le nouvel énoncé porte la date 1832 et montre,
par la manière dont il est écrit, que l'auteur était extrêmement pressé, ce qui
confirme l'assertion que j'ai avancée dans la Note précédente. \hfill \textsc{(A.~Ch.)}}

\begin{center}{\bfseries PROPOSITION IV.}\end{center}

\noindent\textsc{Théorème.} --- \emph{Si l'on adjoint à une équation la valeur numérique
d'une certaine fonction de ses racines, le groupe de
l'équation s'abaissera de manière à n'avoir plus d'autres
permutations que celles par lesquelles cette fonction est invariable.}

En effet, d'après la proposition I, toute fonction connue doit
être invariable par les permutations du groupe de l'équation.

\begin{center}{\bfseries PROPOSITION V.}\end{center}

\noindent\textsc{Problème.} --- \emph{Dans quel cas une équation est-elle soluble par
de simples radicaux?}

J'observerai d'abord que, pour résoudre une équation, il faut
successivement abaisser son groupe jusqu'à ne contenir plus
qu'une seule permutation. Car, quand une équation est résolue,
une fonction quelconque de ses racines est connue, même quand
elle n'est invariable par aucune permutation.

Cela posé, cherchons à quelle condition doit satisfaire le groupe
d'une équation, pour qu'il puisse s'abaisser ainsi par l'adjonction
de quantités radicales.

Suivons la marche des opérations possibles dans cette solution,
en considérant comme opérations distinctes l'extraction de chaque
racine de degré premier.
% GAL1897_W09 stage 2: source-faithful transcription, PDF 072--079 / L0071--L0078.
% Continues Proposition V without reopening the verified PDF 062--071 checkpoint.
